\section{Alcance y metodología}
\label{sec:alcance}

\subsection{Objetivo}

El objetivo general fue evaluar la seguridad del prototipo SOS Clinic, documentar
cada hallazgo con evidencia reproducible y proponer una remediación verificable.
Los objetivos específicos fueron:

\begin{itemize}
  \item Identificar secretos, configuraciones inseguras y vulnerabilidades en las
  dependencias de los siete módulos Java.
  \item Verificar la autenticación, la autorización por rol y el control de acceso
  a objetos en las APIs de negocio.
  \item Evaluar el frontend estático frente a XSS, control de navegación y
  almacenamiento de sesión.
  \item Revisar la exposición de la infraestructura local y el comportamiento del
  agente de IA.
  \item Consolidar los resultados en una matriz única de hallazgos y un plan de
  remediación priorizado.
\end{itemize}

\subsection{Alcance}

\begin{table}[htbp]
\centering
\caption{Alcance incluido y excluido de la evaluación.}
\label{tab:alcance}
\begin{tabular}{@{}L{7.6cm}L{7.6cm}@{}}
\toprule
\textbf{Incluido} & \textbf{Excluido}\\
\midrule
\texttt{ms-gateway}, \texttt{ms-seguridad}, \texttt{ms-pacientes},
\texttt{ms-triage}, \texttt{ms-analitica} y \texttt{ms-ia-agente}. &
\texttt{legacy/} (prototipo Java EE fuera del arranque actual).\\
Frontend estático: sesión, navegación, XSS y consumo de APIs. &
Datos clínicos reales y cualquier entorno de producción.\\
Infraestructura local: PostgreSQL, pgAdmin, Eureka, Mailpit y puertos publicados. &
Redes externas al equipo local y servicios de terceros distintos de Groq.\\
Datos sintéticos y usuarios de prueba creados para la evaluación. &
Ingeniería social, pruebas físicas y denegación de servicio sostenida.\\
\bottomrule
\end{tabular}
\end{table}

La autorización fue otorgada por el responsable del repositorio para ejecutar las
pruebas sobre su entorno local. No se realizaron pruebas contra producción, redes
de terceros ni correo real. El entorno incluye los seis microservicios, el gateway,
el frontend, PostgreSQL, pgAdmin, Eureka y Mailpit; los puertos y las relaciones se
resumen en la figura~\ref{fig:arquitectura}.

\subsection{Metodología}

La evaluación siguió la estructura de PTES y se apoyó en marcos de verificación y
catálogos de debilidades. La tabla~\ref{tab:marcos} resume el uso de cada marco.

\begin{table}[htbp]
\centering
\caption{Marcos de referencia utilizados.}
\label{tab:marcos}
\begin{tabular}{@{}L{5.2cm}L{9.8cm}@{}}
\toprule
\textbf{Marco} & \textbf{Uso en la evaluación}\\
\midrule
PTES & Estructura general de las fases de preparación, análisis, explotación y reporte.\\
OWASP WSTG v4.2 & Casos de prueba web: autenticación, sesión, entrada, errores y configuración.\\
OWASP API Security Top 10 2023 & BOLA, autenticación, mass assignment, BFLA y consumo de recursos.\\
OWASP ASVS 4.0 & Verificación de requisitos de autenticación, sesión, control de acceso y validación.\\
ISO/IEC 27002:2022 & Controles de referencia para la lectura académica y la remediación.\\
CVSS 4.0 & Referencia conceptual para la severidad; el informe fuente no incluye vectores formales por hallazgo.\\
STRIDE/DREAD & Modelado de amenazas por componente y priorización cualitativa.\\
CWE & Clasificación de la debilidad subyacente en cada hallazgo.\\
\bottomrule
\end{tabular}
\end{table}

\subsection{Criterios de severidad}

\begin{table}[htbp]
\centering
\caption{Criterios cualitativos usados para asignar severidad.}
\label{tab:criterios}
\begin{tabular}{@{}L{2.3cm}L{12.6cm}@{}}
\toprule
\textbf{Nivel} & \textbf{Criterio}\\
\midrule
\crit{Crítica} & Compromiso directo de confidencialidad, integridad o autenticación; explotación sin credenciales o con acceso mínimo; afecta datos clínicos o permite administración total.\\
\alta{Alta} & Explotación viable con impacto significativo; requiere una condición razonable (sesión, interacción o acceso a una ruta) y puede derivar en fuga o alteración de datos.\\
\media{Media} & Debilidad de configuración, exposición acotada o abuso limitado; contribuye a una cadena de ataque o reduce la defensa en profundidad.\\
\baja{Baja} & Problema de robustez, manejo de errores o control defensivo; impacto limitado por sí solo.\\
\obs{Observación} & Desajuste funcional, comportamiento esperado mal documentado o falso positivo de una herramienta.\\
\bottomrule
\end{tabular}
\end{table}

\subsection{Reglas de evidencia y ética}

\begin{itemize}
  \item Todas las pruebas se ejecutaron contra \texttt{localhost} y con datos
  sintéticos. El correo se capturó en Mailpit; nunca se usó SMTP real.
  \item Las capturas se sanearon antes de versionarse: tokens, contraseñas,
  cédulas, nombres, correos y códigos 2FA fueron omitidos o enmascarados.
  \item Cada evidencia conserva la fecha, la rama y el commit probado. Los
  archivos crudos se guardaron en \texttt{seguridad/resultados/}.
  \item No se publicaron secretos reales. Los valores sensibles se describen por
  su ubicación y efecto, no por su contenido.
  \item Las pruebas se limitaron a demostrar el hallazgo; no se persistieron
  datos reales ni se degradó el servicio de forma sostenida.
\end{itemize}

\subsection{Limitaciones}

\begin{itemize}
  \item El entorno es de desarrollo y los resultados no equivalen a una auditoría
  de producción. Las cabeceras y la configuración del servidor definitivo pueden
  variar.
  \item ZAP se ejecutó en modo \textit{baseline} (pasivo) sobre el frontend; la API
  se probó de forma manual con scripts reproducibles.
  \item El SCA utilizó el caché local de Maven porque Maven Central respondió
  \texttt{429 Too Many Requests} durante la ejecución inicial.
  \item Trivy no detectó \texttt{docker-compose.yml}; la configuración de
  infraestructura se revisó manualmente.
  \item El informe fuente no calcula vectores CVSS individuales. Las correcciones
  posteriores en \texttt{main} se verificaron con \texttt{mvn test}, Trivy,
  gitleaks y pruebas E2E (sección~\ref{sec:estado-remediacion}), pero no
  constituyen una auditoría independiente ni una re-ejecución completa de todos
  los casos.
  \item No se evaluó \texttt{legacy/} ni se realizaron pruebas de denegación de
  servicio sostenidas, ingeniería social o acceso físico.
\end{itemize}

